
\chapter{Платежи по~QR и~кошельки}\label{ch:qr-wallets}
\highlight{QR-код и~платёжное приложение не~образуют отдельной платёжной
  системы. Фактический расчёт выполняет карточная сеть, A2A-перевод или~закрытый
кошелёк}.
Карточные системы остаются основой потребительских расчётов; альтернативы
выигрывают там, где снижают комиссию, дают рассрочку или решают локальную
задачу.

\section{Таксономия методов оплаты}\label{sec:payment-methods-taxonomy}

Метод оплаты задают три признака: кто инициирует списание, где наступает
финальность и~кто несёт риск спора. По~этим признакам карта, банковский
перевод, кошелёк, рассрочка и~цифровые деньги расходятся сильнее, чем по~виду
кнопки или QR-кода на~кассе.

Первый класс~--- карточные платежи: Visa, Mastercard, Мир, UnionPay и~American
Express. Продавец инициирует списание, отправляя запрос по~карточным реквизитам
клиента и~получая либо авторизацию, либо отказ. Преимущества карт~---
глобальная инфраструктура, возвратный платёж (chargeback~--- сетевой возврат
  средств эмитентом по~спорной операции, в~отличие от~коммерческого возврата
продавца; см.~гл.~\ref{ch:disputes}) и~развитая практика PCI~DSS (Payment Card
  Industry Data Security Standard~--- отраслевой стандарт защиты карточных данных;
область применимости разбирается в~гл.~\ref{ch:pci-dss}). Ограничение~--- стоимость
приёма для~продавца.

Второй класс~--- переводы со~счёта на~счёт (account-to-account, A2A), где
клиент подтверждает перевод со~своего счёта на~счёт продавца~--- это
\textit{push}-модель, в~отличие от~\textit{pull} карточных платежей, где
списание по~реквизитам клиента инициирует продавец. К~A2A относятся СБП, SEPA Instant Credit
Transfer, Pix и~UPI~\cite{sbp-main,sepa-instant,pix,upi}. Преимущества A2A~---
финальность (после зачисления перевод нельзя односторонне отменить
по~правилам системы) и~низкие комиссии; ограничение~--- более узкая защита
покупателя.

Третий класс~--- кошельки: надстройка над базовой платёжной инфраструктурой.
\highlight{Закрытый кошелёк (\textit{closed-loop wallet}) учитывает средства
  на~собственном балансе и~проводит расчёт внутри своего реестра, как Alipay
  и~WeChat Pay. Открытый кошелёк (\textit{open-loop wallet}) заменяет
реквизиты карты платёжным токеном}. По~этой модели работают Apple Pay и~Google Pay;
о~платёжной токенизации, на~которой строится открытый кошелёк, подробно
см.~гл.~\ref{ch:tokenization}. Для российских карт после марта 2022~года оба
сервиса недоступны; их функции частично заменяют Mir Pay и~SBPay
(см.~раздел~\ref{sec:russia-wallets}).

Четвёртый класс~--- BNPL (Buy Now Pay Later): провайдер оплачивает продавцу всю
сумму сразу, а~клиент возвращает её частями по~собственному графику. Для
продавца BNPL похож на~эквайринг с~иной комиссией, но~экономическая модель
другая: ключевой элемент~--- кредитное решение провайдера, подробно разобранное
в~главе~\ref{ch:bnpl}.

Пятый класс~--- наличный сценарий (\textit{cash-based}), в~котором платёж
начинается онлайн, но~завершается офлайн через ваучер или платёжный
идентификатор.
Типичный пример~--- OXXO в~Мексике, где покупатель получает код платежа
и~завершает операцию в~физической точке~\cite{oxxo-pay}. Подтверждение приходит
продавцу не~сразу, а~после того, как покупатель внесёт деньги в~кассе.

Шестой класс~--- цифровые деньги вне~банковского депозита; он объединяет два
разных явления. Цифровая валюта центрального банка (central bank digital
currency, CBDC)~--- цифровая форма обязательства центрального банка, и~цифровой
рубль разобран в~главе~\ref{ch:cbdc}. Криптовалюты и~стейблкоины строятся
на~иной модели эмиссии и~регулирования. Общий признак с~CBDC здесь один:
такие инструменты не~сводятся к~обычному банковскому счёту.

\begin{table}[H]
  \caption{Сравнение методов оплаты по~ключевым параметрам}\label{tab:payment-methods-matrix}
  \begingroup
  \tablesetup
  \scriptsize
  \setlength{\tabcolsep}{2pt}
  \begin{minipage}{\textwidth}
    \begin{tabularx}{\textwidth}{
        @{}B{0.170\textwidth}
        P{0.075\textwidth}
        P{0.095\textwidth}
        P{0.090\textwidth}
        P{0.085\textwidth}
        P{0.145\textwidth}
        P{0.100\textwidth}
        P{0.155\textwidth}@{}
      }
      \toprule
      \headrow
      \thd{Метод оплаты} & \thd{Тип} & \thd{Финальность} & \thd{Возвратный платёж} & \thd{PCI
      scope} & \thd{Международность} & \thd{Скорость} & \thd{Комиссия ТСП} \\
      \midrule
      Карты &
      \cellcolor{Bad!10}pull &
      \cellcolor{Bad!10}нет &
      \cellcolor{Bad!10}есть &
      \cellcolor{Bad!10}да &
      \cellcolor{Good!14}глобальный &
      \cellcolor{Good!14}мс &
      \cellcolor{Bad!10}высокая \\
      \rowsep
      A2A (СБП/Pix) &
      \cellcolor{Good!14}push &
      \cellcolor{Good!14}да &
      \cellcolor{Good!14}нет &
      \cellcolor{Good!14}нет &
      \cellcolor{Ink!6}региональный &
      \cellcolor{Good!14}сек &
      \cellcolor{Good!14}низкая \\
      \rowsep
      Closed-loop кошелёк &
      \cellcolor{Ink!6}push/pull &
      \cellcolor{Ink!6}зависит &
      \cellcolor{Ink!6}ограничен &
      \cellcolor{Good!14}нет &
      \cellcolor{Bad!10}нишевый &
      \cellcolor{Good!14}мс &
      \cellcolor{Ink!6}средняя \\
      \rowsep
      BNPL &
      \cellcolor{Bad!10}pull &
      \cellcolor{Bad!10}нет &
      \cellcolor{Bad!10}есть &
      \cellcolor{Ink!6}нет &
      \cellcolor{Ink!6}региональный &
      \cellcolor{Good!14}сек &
      \cellcolor{Ink!6}средняя \\
      \rowsep
      Cash-based &
      \cellcolor{Good!14}push &
      \cellcolor{Good!14}да &
      \cellcolor{Good!14}нет &
      \cellcolor{Good!14}нет &
      \cellcolor{Bad!10}нишевый &
      \cellcolor{Bad!10}часы / дни &
      \cellcolor{Ink!6}нет данных \\
      \rowsep
      CBDC (цифровой рубль) &
      \cellcolor{Good!14}push &
      \cellcolor{Good!14}да &
      \cellcolor{Good!14}нет &
      \cellcolor{Good!14}нет &
      \cellcolor{Bad!10}региональный &
      \cellcolor{Ink!6}сек / мин &
      \cellcolor{Ink!6}нет данных \\
      \rowsep
      Crypto &
      \cellcolor{Good!14}push &
      \cellcolor{Good!14}да &
      \cellcolor{Good!14}нет &
      \cellcolor{Good!14}нет &
      \cellcolor{Ink!6}глобальный / нишевый &
      \cellcolor{Ink!6}мин / часы &
      \cellcolor{Ink!6}нет данных \\
      \bottomrule
    \end{tabularx}
    \par\smallskip
    \tablenote{Зелёный цвет показывает сравнительное преимущество,
      красный указывает на~ограничение, серый~--- на~нейтральную или
      зависящую от~провайдера зону. Колонки: \enquote{Тип}~---
      \textit{push} (инициирует и~подтверждает клиент) или \textit{pull}
      (списание инициирует продавец); \enquote{Финальность}~---
      безотзывность перевода после зачисления по~правилам системы;
      \enquote{Возвратный платёж}~--- наличие сетевого механизма
      \textit{chargeback}, а~не коммерческого возврата
      (см.~гл.~\ref{ch:disputes}); \enquote{PCI~scope}~--- попадание
      продавца в~область применимости PCI~DSS (см.~гл.~\ref{ch:pci-dss}).
      Приведённые оценки меняются в~зависимости от~конкретного
    провайдера, маршрута и~географии.}
  \end{minipage}\par
  \endgroup
\end{table}

\section{Платежи по~QR: три модели}\label{sec:qr-models}

Фраза \enquote{мы принимаем QR} без~указания инфраструктуры не~сообщает, кто
подтверждает платёж, где наступает финальность и~кто разбирает спорные
операции. \highlight{QR-код только передаёт платёжные данные без~оборудования
  NFC или~ручного ввода реквизитов; расчёт
выполняет A2A-перевод, закрытый кошелёк или~карточная транзакция}.

QR-код в~платежах используется в~двух сценариях, которые отличаются тем, кто
инициирует сканирование. При~\textit{merchant-presented mode} (MPM) продавец
показывает код, а~клиент сканирует его камерой телефона.
При~\textit{customer-presented mode} (CPM) клиент открывает динамический код
в~приложении, а~продавец считывает его сканером на~кассе. Термины MPM и~CPM
закреплены в~спецификациях QR-кодов EMVCo~\cite{emvco-qr-mpm,emvco-qr-cpm}. MPM
дешевле в~развёртывании, потому что достаточно напечатанной наклейки или экрана
на~кассе, тогда как CPM быстрее на~высоком потоке покупателей, но~требует
от~продавца оборудования для~считывания. Обе модели встречаются в~каждой
из~трёх схем ниже, но~преобладают по-разному.

\subsection*{СБП: QR поверх A2A-перевода}

В~СБП QR-код ведёт в~мгновенный перевод со~счёта клиента на~счёт ТСП через
инфраструктуру НСПК~\cite{sbp-main}. Система работает как \textit{push}-платёж:
клиент инициирует и~подтверждает списание.

Статический QR~--- напечатанный код, который содержит идентификатор ТСП,
но~не~содержит сумму. Клиент сканирует код в~приложении своего банка, вводит
сумму вручную и~подтверждает перевод. Такой код подходит малому бизнесу:
кассовое ПО не~обязательно, достаточно наклейки. Динамический QR
генерируется кассовой системой под конкретную покупку и~содержит реквизиты
получателя, сумму и~назначение платежа. Клиент сканирует код, видит
предзаполненные данные и~подтверждает одним действием, а~кассовая система
получает обратный вызов об~успешном переводе за~несколько
секунд~\cite{sbp-faq-business}.

Помимо QR, СБП поддерживает оплату по~платёжной ссылке, по~NFC и~по~привязке
счёта в~приложении SBPay~\cite{sbpay-app,sbp-main}. Все эти каналы ведут
к~одному и~тому же A2A-переводу внутри инфраструктуры НСПК; различается только
способ, которым клиент передаёт согласие на~списание.

Тарифы C2B-переводов (consumer-to-business~--- от~физического лица в~пользу
ТСП) в~СБП устанавливает Совет директоров Банка России. Ставка межбанковского
вознаграждения зависит от~MCC (Merchant Category Code, код категории торговой
точки) ТСП и~составляет от~0\,\% до~0{,}7\,\% от~суммы перевода.
Для платежей в~пользу государства она равна 0\,\%, для услуг ЖКХ ограничена
0{,}2\,\%, для социально значимых категорий (продукты, лекарства, медицина
и~т.\,п.)~--- 0{,}4\,\%, для прочих категорий~---
0{,}7\,\%~\cite{cbr-sbp-c2b-tariffs}. Для сравнения, \textit{interchange}
(межбанковская комиссия в~карточной сети,
см.~раздел~\ref{sec:ecosystem-interchange}) в~карточном эквайринге начинается
от~1\,\% и~превышает 2\,\%.

Механизма возвратного платежа в~СБП нет, поскольку перевод финален с~момента
зачисления на~счёт получателя. При возврате покупки продавец инициирует
обратный перевод через свой банк~\cite{sbp-faq-business}.
Ошибочный перевод не~тому получателю также требует согласия получателя
на~возврат, если иной порядок не~установлен законом или решением
суда~\cite{cbr-sbp-faq-mistaken-transfer}. Эта модель выгоднее продавцу,
но~смещает риск спора на~покупателя.

\subsection*{Alipay и~WeChat Pay: QR поверх закрытого
кошелька}\label{sec:qr-wallet-closed-loop}

В~Alipay и~WeChat Pay QR-код ведёт не~в~банковский перевод, а~в~платёж внутри
закрытой системы кошелька. Средства пользователя учитываются на~балансе
кошелька, который пополняется с~банковского счёта или карты; оператор списывает
сумму без~выхода в~межбанковский клиринг.

Оба сценария, MPM и~CPM, здесь используются, причём в~документации
Alipay+ merchant-presented mode описан как основной сценарий для сторонних ТСП.
Продавец показывает QR-код, клиент сканирует его в~приложении кошелька,
оператор списывает средства с~баланса клиента и~зачисляет их на~счёт
ТСП~\cite{alipayplus-mpm}. Customer-presented mode распространён
в~высокопоточных офлайн-точках вроде сетевых магазинов, автоматов
и~общественного транспорта. Клиент показывает динамический штрихкод
из~приложения, кассир сканирует его, и~оплата проходит без~дополнительного
подтверждения при сумме ниже порога беспарольного списания (для WeChat Pay~---
  порядка нескольких тысяч юаней с~дневным лимитом на~число операций, для Alipay+
оператор задаёт порог с~учётом валюты расчёта).

Оператор кошелька совмещает роли, которые в~карточной сети разделены между
эмитентом и~эквайером. Он учитывает балансы клиентов и~ТСП, проводит расчёт
внутри своего реестра и~сам определяет правила диспутов. Для клиента путь
от~пополнения до~возврата остаётся внутри одного приложения; продавец
интегрируется через API оператора кошелька, минуя эквайера.

Трансграничное расширение устроено как партнёрская сеть, в~которой Alipay+
объединяет локальные кошельки (GCash, TrueMoney, KakaoPay и~другие), а~их
пользователи платят в~точках, принимающих Alipay, за~пределами домашней страны.
Продавец подключается к~одному контракту с~локальным партнёром Alipay+
и~получает доступ к~пользователям нескольких кошельков~\cite{alipayplus-mpm}.
Модель отличается от~карточной модели Visa или Mastercard: вместо единого
глобального протокола работает федерация кошельков, каждый из~которых действует
в~собственной расчётной инфраструктуре.

\subsection*{Pix: QR, встроенный в~национальную платёжную
инфраструктуру}\label{sec:qr-pix}

Pix как платёжная схема разобран в~разделе~\ref{sec:pix-system}. В~отличие
от~FedNow и~SCT Inst, где правила схемы заканчиваются на~pacs.008, в~Pix
QR-код~--- штатный интерфейс системы: формат BR Code задаёт сам
BCB~\cite{pix-manual-iniciacao}.

BCB определяет два типа QR. Статический многоразовый код содержит идентификатор
получателя и~опционально сумму. Динамический одноразовый генерируется
под~конкретную транзакцию с~полным набором данных; после исполнения кассовая
система получает обратный вызов. Помимо сканирования, Pix поддерживает Pix
Copia e Cola~--- текстовое представление тех же данных, которое копируется
и~вставляется в~банковское приложение~\cite{pix-faq}. Этот вариант нужен, когда
приложение плательщика и~источник кода находятся на~одном экране и~камерой
сканировать нечего.

Комиссия для физических лиц за~переводы и~платежи через Pix равна нулю
по~решению регулятора. Для юридических лиц тарифы определяются
банком-участником, но~ниже карточного эквайринга. Возвратный платёж
как сетевая процедура отсутствует~--- Pix работает как \textit{push}-платёж
и~перевод финален с~момента зачисления на~счёт получателя; возврат по~решению
противомошеннической проверки выполняется через отдельный MED-канал схемы
(см.~раздел~\ref{subsec:pix-med}).

\subsection*{Сравнение трёх моделей}

Один и~тот же носитель, QR-код, ведёт к~разным расчётным моделям
(табл.~\ref{tab:qr-three-models}). СБП и~Pix проводят платёж через
межбанковскую инфраструктуру, управляемую центральным банком, тогда как Alipay
или WeChat Pay замыкают расчёт внутри собственного реестра. A2A-модель даёт
мгновенный \textit{push}-перевод с~ограниченной регулятором комиссией
и~без~возвратного платежа; модель закрытого кошелька даёт оператору контроль
над правилами диспутов, ценообразованием и~международным расширением.

\begin{table}[htbp]
  \caption{Три модели платежей по~QR}\label{tab:qr-three-models}
  \begingroup
  \tablesetup
  \setlength{\tabcolsep}{3pt}
  \begin{minipage}{\textwidth}
    \begin{tabularx}{\textwidth}{
        @{}B{0.16\textwidth}
        P{0.255\textwidth}
        P{0.255\textwidth}
        P{0.255\textwidth}@{}
      }
      \toprule
      \headrow
      \thd{\thd{}}
      & \thd{СБП} & \thd{Alipay / WeChat Pay} & \thd{Pix} \\
      \midrule
      Инфраструктура & A2A-перевод (межбанковский) & Закрытый кошелёк & A2A-перевод (межбанковский) \\
      \rowsep
      Оператор & НСПК (Банк России) & Ant Group / Tencent & Banco Central do Brasil \\
      \rowsep
      Расчёт & Через НСПК, секунды & Внутри реестра оператора & SPI при BCB, секунды \\
      \rowsep
      Финальность & Мгновенная (push) & Мгновенная (внутри системы) & Мгновенная (push) \\
      \rowsep
      Возвратный платёж & Нет & Правила оператора & Нет \\
      \rowsep
      Статический QR & Да & Да & Да \\
      \rowsep
      Динамический QR & Да & Да & Да \\
      \rowsep
      Основной сценарий QR & MPM & MPM + CPM & MPM \\
      \rowsep
      Комиссия C2B & 0\,\%--0{,}7\,\% (регулятор) & Устанавливает оператор & 0\,\% (физ.\ лица) \\
      \rowsep
      Трансграничность & Региональный & Федерация кошельков & Региональный \\
      \bottomrule
    \end{tabularx}
  \end{minipage}\par
  \endgroup
\end{table}

\section{Платёжные приложения и~кошельки в~России}\label{sec:russia-wallets}

После марта 2022~года в~России сложилась собственная структура платёжных
приложений и~кошельков. Она заменяет часть функций ушедших Apple Pay и~Google
Pay и~отражает специфику внутрироссийской инфраструктуры с~картами
\enquote{Мир}, СБП и~цифровым рублём. Набор поддерживаемых функций и~состав
банков-партнёров у~сервисов 2024--2026~годов меняется от~релиза к~релизу.

\textbf{Mir Pay}~--- кошелёк NFC на~Android поверх токенизации НСПК
(см.~раздел~\ref{sec:mir-pay}). \textbf{SBPay}~--- мобильное приложение НСПК
поверх СБП с~оплатой по~QR, табличке NFC, кнопке и~ссылке за~счёт привязанного
счёта, доступен и~на~Android, и~на~iOS (см.~раздел~\ref{sec:sbp-sbpay}).

\textbf{SberPay}~--- кошелёк и~платёжный сервис ПАО \enquote{Сбербанк}, который
поддерживает оплату на~сайтах партнёров через кнопку SberPay и~в~подключённых
офлайн-точках. Для держателей карт \enquote{Мир} Сбербанка также реализованы
сценарии оплаты по~QR-коду и~NFC через мобильное приложение
банка~\cite{sber-pay}.

\textbf{T-Pay}, бывший \enquote{Tinkoff Pay},~--- платёжный сервис АО
\enquote{ТБанк}, бывшего Тинькофф Банка. Сервис работает на~сайтах партнёров
и~в~офлайн-точках, в~том числе с~использованием NFC на~Android, а~сценарии
оплаты опираются на~карты и~счета клиентов банка~\cite{tpay}.

\textbf{AlfaPay}~--- платёжный сервис АО \enquote{Альфа-Банк} для~карт банка
ПС~Мир. На~An\-droid~--- NFC через HCE в~приложении банка с~пулом
одноразовых ключей: банк заявляет до~10~транзакций между сессиями онлайн
(LUK-пул разобран
в~разделе~\ref{sec:contactless-mobile})~\cite{alfa-pay-offline-2025}.
На~iPhone~--- BLE поверх \enquote{Вжуха} Сбера. Если в~смартфоне нет NFC,
доступен стикер Альфа-банка~--- NFC-наклейка как отдельная дебетовая карта
ПС~Мир; он работает и~на~iPhone, и~на~любом телефоне без~NFC. Для
ТСП AlfaPay~--- кнопка в~интернет-эквайринге Альфа-Банка: REST-вызов
на~платёжный шлюз открывает приложение банка по~\textit{deep-link},
и~покупатель подтверждает оплату в~приложении банка, не~вводя реквизиты карты
на~сайте ТСП~\cite{alfa-pay-web-api-2025}.

\textbf{Yandex Pay}~--- платёжный сервис ООО \enquote{Яндекс.Пэй}, встраиваемый
в~форму оплаты партнёрских сайтов и~ориентированный на~оплату ранее
сохранёнными картами без~повторного ввода реквизитов (\textit{frictionless}).
С~2024~года он~интегрирован с~сервисом рассрочки \enquote{Сплит}
(раздел~\ref{sec:bnpl-russia})~\cite{yandex-pay,yandex-split}.

\textbf{ЮKassa}~--- платёжный сервис ООО НКО \enquote{ЮМани}, дочерней
структуры Сбербанка. Работает как PSP для~бизнеса с~приёмом карт, платежей СБП
и~BNPL-сценариев. Часть инфраструктуры общая с~одноимённым клиентским кошельком
\textbf{ЮMoney}, который остаётся электронным средством платежа для розничных
пользователей и~наследует Яндекс.Деньги~\cite{yookassa,yoomoney}.

\textbf{Биоэквайринг и~оплата по~лицу}~--- сценарий, который НСПК и~Банк России
развивают на~основе Единой биометрической системы (ЕБС). Покупатель оплачивает
покупку, подтверждая личность биометрией на~терминале, без~карты
и~телефона~\cite{nspk-biopay-overview,ebs-overview}. К~2025~году платформа НСПК
объединяет несколько биоэквайреров и~эмитентов, оплата работает в~кассах
самообслуживания X5 (\enquote{Перекрёсток}, \enquote{Пятёрочка})
и~на~турникетах московского метро. Оборот канала на~порядки ниже, чем у~карт
и~СБП, но~он уже используется в~рознице и~транспорте.

Для ТСП в~России в~2024--2026~годах практический базовый набор~--- приём карт
\enquote{Мир} через эквайринг, СБП по~QR или~кнопке, один или несколько
кошельков (SberPay, T-Pay, Yandex Pay) и~BNPL для подходящих категорий;
на~отдельных типах ТСП добавляется биоэквайринг.
Apple Pay и~Google Pay для российских карт в~этот набор больше не~входят.

\section{Структура способов оплаты в~России и~СНГ}\label{sec:payments-mix-cis}

Россия сочетает массовую карточную инфраструктуру и~растущие альтернативы.
Банк России указывает, что карты остаются самым популярным
средством безналичной оплаты, а~платежи по~QR-коду в~2024~году достигли
4{,}1~трлн~рублей при~21{,}6~млрд~рублей оплат
по~биометрии~\cite{cbr-results-2024-payments}. В~Казахстане национальная
статистика отдельно выделяет крупный поток безналичных платежей по~QR наряду
с~карточными операциями~\cite{nbk-payment-cards-2025jan}.
Узбекистан опирается на~две национальные карточные системы~--- HUMO
и~UzCard~\cite{humo-faq,uzcard-about}. Армения сохраняет собственную карточную
инфраструктуру ArCa, и~оператор системы указывает, что банки-участники
выпускают карты ArCa наряду с~Visa, Mastercard и~другими международными
системами~\cite{arca-about}.

При локализации продукта в~регионе структуру способов оплаты нужно оценивать
отдельно по~каждой стране: российский набор нельзя переносить на~Казахстан,
а~армянскую модель на~Узбекистан без~проверки локальной инфраструктуры, роли QR
и~регуляторных требований.

\practicenote{%
  При локализации в~новой стране СНГ проверьте распределение
  BIN трафика, наличие массового локального A2A-метода
  и~требования к~обязательному приёму локальных карт. Эти три фактора
  определяют экономику приёма и~конверсию сильнее остальных.

}

Правовой статус локальной карточной системы задают три требования:
обязательной маршрутизации всех внутренних карточных операций через
национальный коммутатор (как в~России~--- через НСПК), обязательного приёма
национальных карт всеми ТСП определённого размера и~отдельных тарифов
на~C2B-переводы по~QR, утверждаемых регулятором. В~Казахстане локальный
коммутатор~--- ISPC АО~\enquote{Национальная платёжная корпорация Национального
банка Республики Казахстан}; он работает с~30~июня 2022~года и~обрабатывает
внутриказахстанские карточные операции~\cite{nbk-ispc-launch-2022}.
В~Узбекистане параллельно работают две
национальные карточные системы~--- HUMO (Национальный межбанковский
процессинговый центр) и~UzCard (Единый общереспубликанский процессинговый
центр)~\cite{humo-faq,uzcard-about}; правила локальной маршрутизации задают
Национальный банк Казахстана и~Центральный банк Узбекистана отдельно. В~Армении
ArCa~--- единый процессинг для собственного бренда и~для Visa, Mastercard,
American Express, JCB и~\enquote{Мир}; банки-участники эмитируют ArCa
параллельно с~международными системами~\cite{arca-about}. Доля карт ArCa
в~эмиссии (около 16\,\% на~01.10.2025) ниже доли карт \enquote{Мир}
в~России (около 62\,\% эмиссии и~около 64\,\% стоимости внутрироссийских
операций по~картам по~данным Банка России~\cite{cbr-onrnps-2025}).

\section{Почему карты остаются основным инструментом}\label{sec:card-dominance}

Карты сохраняют центральное место в~потребительских платежах из-за четырёх
инерций.

Сетевые эффекты: глобальные карточные сети уже выстроили массовую эмиссию,
приём и~привычный пользовательский сценарий. Продавец, который принимает карты,
доступен большинству покупателей, а~новый метод должен одновременно получить
поддержку и~у~эквайеров, и~у~эмитентов.

Международность: карта, выпущенная в~одной стране, работает во~множестве других
по~единому протоколу, тогда как локальные A2A-системы вроде СБП, Pix и~UPI
не~взаимозаменяемы.

Защита покупателя через возвратный платёж: механизм, который A2A-системы
не~воспроизводят без~изменения самой идеи финальности. Покупатель, оплативший
картой незнакомому продавцу, имеет формальный механизм возврата. Покупатель,
оплативший через A2A, может требовать возврат по~договору или специальным
правилам схемы, но~возвратного платежа уровня карточной системы здесь нет.

Инфраструктурная инерция: мировая сеть POS-терминалов (Point of Sale, кассовый
терминал в~физической точке обслуживания) уже выстроена под карточные сети.
Переоснащение или добавление нового метода окупается при массовом спросе,
а~массовый спрос не~возникает, пока продавцы не~видят достаточного
потока покупателей. Разорвать эту зависимость помогает регулирование или
встроенность в~уже массовый кошелёк.

Альтернативные методы выигрывают в~нишах: A2A снижает комиссию, наличные
сценарии обслуживают клиентов без~банковского счёта, CBDC добавляет условные
платежи, BNPL даёт рассрочку, а~криптоактивы поддерживают программируемые
международные сценарии. Глобальную универсальность и~инфраструктурный охват
карт они пока не~воспроизводят, поэтому работают рядом с~картами там, где их
преимущество перевешивает универсальность карточного приёма.

Выбор метода привязан к~рынку и~чеку. Внутри России при среднем чеке ниже
100\,000~руб. СБП даёт комиссию от~0 до~0{,}7\,\% в~зависимости от~категории ТСП
и~мгновенное зачисление, а~карты остаются резервным методом для клиентов
без~мобильного банка. На~международных рынках карты остаются основным методом,
а~A2A-системы других стран, например Pix или UPI в~Индии, требуют отдельной
интеграции на~каждый рынок. При среднем чеке 10\,000--150\,000~руб.
и~критичной для~продавца конверсии BNPL повышает конверсию за~счёт рассрочки,
но~добавляет кредитный риск на~стороне провайдера и~задержку в~расчёте. Для
покупателей без~банковского счёта, например в~части стран Азии и~Африки, нужны
наличные через агентскую сеть или мобильные деньги. Карты остаются базовым
методом, а~остальные добавляются там, где комиссия, рассрочка, доступ без~счёта
или финальность дают измеримый выигрыш.
